% Chapter 1 exercises. The answers are in answers/barcs_answers_ch1.pdf.
\section{Exercises for Chapter 1}

These exercises use a new screen. Work on paper. Show each calculation, and
write one or two sentences when a question asks you to explain. The answers
are in a separate answer sheet.

The tags tell you what to do: \tagcalc{} work out a number,
\tagexplain{} say why in your own words, \tagwrite{} write down a table or
a matrix, \tagdecide{} choose and give a reason.

\begin{center}
\fbox{\begin{minipage}{0.92\linewidth}
\textbf{The screen.} A lab tests a new drug. They grow cells with a guide
library, split them into four flasks, and treat two flasks with the drug.
After two weeks they read the guides. Here are the results for one guide,
\textbf{guide g7}:

\smallskip
\centering
\begin{tabular}{lcrr}
\toprule
Library & Group & Guide count $K$ & Total count $N$\\
\midrule
1 & control & 150 & 1{,}000{,}000\\
2 & control & 210 & 1{,}400{,}000\\
3 & drug    &  60 &   800{,}000\\
4 & drug    &  45 &   600{,}000\\
\bottomrule
\end{tabular}
\end{minipage}}
\end{center}

\subsection*{Part A: From counts to a change}
\begin{exlist}[series=ex1]
\item \tagcalc
\begin{parts}
\item Calculate the proportion of guide g7 in each library, in reads per
  million.
\item Calculate the combined proportion of each group (add the counts, add
  the totals, then divide).
\item Calculate the fold change (drug $\div$ control) and the $\log_2$
  fold change.
\end{parts}

\item \tagexplain{} A student looks only at the counts. Library~2 has 210
reads and library~3 has 60. The student says: ``The drug cut guide g7 by
$210 \div 60 = 3.5$ times.'' What is wrong with this? Why must we divide by
$N$ before we compare libraries?

\item \tagcalc{} \tagexplain
\begin{parts}
\item Calculate the log-odds of the control proportion,
  $\ln\bigl(0.00015 / (1 - 0.00015)\bigr)$.
\item BARCS works with log-odds, not with proportions. Give two reasons.
  (Hint: what values can a proportion have, and what values can a log-odds
  have? What happens to ``half'' and ``double'' on the log scale?)
\end{parts}
\end{exlist}

\subsection*{Part B: Noise}
\begin{exlist}[resume*=ex1]
\item \tagcalc{} \tagexplain{} Four repeated counts are 140, 150, 160 and
150. Their average is 150.
\begin{parts}
\item Calculate the distance of each count from 150, and add the four
  distances.
\item Calculate the average squared distance (the variance) and its square
  root (the standard deviation).
\item Use part (a) to explain why the variance uses \emph{squared}
  distances.
\end{parts}

\item \tagcalc{} Use library~1: $N = 1{,}000{,}000$ and the expected
proportion $\mu = 0.00015$.
\begin{parts}
\item Calculate the counting noise $N \times \mu \times (1-\mu)$ and its
  standard deviation.
\item BARCS finds $\rho = 0.000001$. Calculate the extra factor
  $1 + (N-1) \times \rho$, the total variance, and the standard deviation.
\item Calculate the biological part
  $N \times (N-1) \times \rho \times \mu \times (1-\mu)$ on its own. Check
  that counting part $+$ biological part $=$ total.
\end{parts}

\item \tagexplain
\begin{parts}
\item In words, what does $\rho = 0.000001$ say about two reads from the
  same library?
\item What would $\rho$ be if every flask had exactly the same true
  proportion of guide g7?
\end{parts}

\item \tagcalc{} \tagexplain{} A tiny library has $N = 4$ reads, called
A, B, C and D.
\begin{parts}
\item List every ordered pair of two different reads (A--B and B--A are
  different pairs). How many are there?
\item Check your count with the formula $N \times (N-1)$.
\item Why is the formula $N \times (N-1)$ and not $N \times N$?
\end{parts}

\item \tagdecide{} The lab can pay for 8 million reads in each group. Keep
$\mu = 0.00015$ and $\rho = 0.000001$. They have three options:
\begin{center}
\begin{tabular}{lcc}
\toprule
Option & Replicates per group & Reads per library\\
\midrule
A & 2 & 4{,}000{,}000\\
B & 4 & 2{,}000{,}000\\
C & 8 & 1{,}000{,}000\\
\bottomrule
\end{tabular}
\end{center}
\begin{parts}
\item For one library, the standard deviation of the proportion is
  $\sqrt{\text{variance of } K} \div N$. Calculate it (in reads per million)
  for $N = 4$, 2 and 1 million.
\item The average of $n$ replicates has the standard deviation of one
  library divided by $\sqrt{n}$. Calculate it for A, B and C.
\item Which option measures the group proportion most precisely? Explain
  why, using the extra factor $1 + (N-1) \times \rho$.
\end{parts}
\end{exlist}

\subsection*{Part C: The design matrix}
\begin{exlist}[resume*=ex1]
\item \tagwrite{} A second experiment has 3 control libraries and 3 drug
libraries. Write the design matrix $X$ with the columns ``intercept'' and
``drug''. How many degrees of freedom does the model have?

\item \tagwrite{} \tagexplain{} A third experiment uses cells from two
donors. Each donor gives one control library and one drug library.
\begin{parts}
\item Write $X$ with the columns ``intercept'', ``donor 2'' and ``drug''.
  Put the libraries in this order: donor~1 control, donor~1 drug, donor~2
  control, donor~2 drug.
\item What is the job of the ``donor 2'' column? What could go wrong
  without it?
\item How many degrees of freedom are left? What does this mean for the
  test?
\end{parts}

\item \tagcalc{} For guide g7, BARCS finds $\beta_0 = -8.8048$ and
$\beta_1 = -0.6931$.
\begin{parts}
\item Use the drug row of $X$, $(1, 1)$, to calculate the predicted
  log-odds $\eta$ of a drug library.
\item Change $\eta$ into the proportion $\mu = e^{\eta} / (1 + e^{\eta})$.
  Give it in reads per million.
\item Calculate the expected count of library~3 ($N = 800{,}000$). Compare
  it with the real count.
\end{parts}

\item \tagexplain{} Calculate $e^{\beta_1}$. What does this number mean for
the experiment? Why is it useful that $\beta_1$ is on the log-odds scale?
\end{exlist}
